\documentclass[10pt,a4paper]{amsart}
\usepackage[margin=19mm]{geometry}
\usepackage{amsmath,amssymb,mathtools}
\usepackage[scr=rsfs,cal=euler]{mathalfa}
\usepackage{xfrac}
\usepackage[unicode,hidelinks,bookmarks=false]{hyperref}
\newcommand{\ie}{i.e.\ }
\newcommand{\cf}{cf.\ }
\newcommand{\resp}{respectively\ }
\newcommand{\Subsection}{\subsection}
\providecommand{\nobreakdash}{\nobreak\hbox{-}}
\setlength{\emergencystretch}{2em}
\multlinegap0pt
\usepackage{bm}
\usepackage{bbm}
\usepackage{dsfont}
\usepackage{xspace}
\usepackage{cancel}
\usepackage{mathtools}
\usepackage{xcolor}
\usepackage{amsthm}
\makeatletter
\@ifundefined{theorem}{\newtheorem{theorem}{Theorem}}{}
\@ifundefined{definition}{\newtheorem{definition}[theorem]{Definition}}{}
\makeatother
\title[On modules over a Hopf brace]{On modules over a Hopf brace}
\author{Ram\'on Gonz\'alez Rodr\'iguez, Brais Ramos P\'erez, Ana Bel\'en Rodr\'iguez Raposo}
\date{Source: arXiv:2603.22932v1 (CC BY 4.0)}
\begin{document}
\maketitle

\noindent\emph{Source and licence.} On modules over a Hopf brace. Version arXiv:2603.22932v1 (CC BY 4.0), distributed under the Creative Commons Attribution 4.0 International licence, as recorded in the arXiv licence metadata for this exact version and shown on the abstract page https://arxiv.org/abs/2603.22932v1. Full rights notice: licenses/arxiv-2603.22932-CC-BY-4.0.txt.\par\smallskip

\noindent\textit{Editorial scope.} Counted unit: the Theorem whose source label is \texttt{usualtensoraction} in the article (source0.tex lines 731--740, source label \texttt{usualtensoraction}), delivered together with the article's own complete original proof of it (source0.tex lines 741--815, one \texttt{proof} environment of 75 lines). No mathematics is written, repaired or re-derived: statement and proof are the article's own text line for line. Required context retained uncounted: a-antip2 (lines 281--284); compatmodH (lines 621--623); GammaMdef (lines 625--627); ccclassGammaM (lines 721--730). Every definition, display and label of those lines is retained unchanged. Omitted from this package and not redistributed: 1--280, 285--620, 624--624, 628--720, 816--1101. Nothing inside a retained excerpt is omitted or altered: every retained body is delivered exactly as the article's lines are written, so no omission receipt of any kind accompanies this package. Citations: the retained excerpts carry no citation command at all, so no bibliography entry of the article is transcribed and no reference is redistributed. Reference adaptation: none was needed and none was made; every reference inside the delivered unit resolves inside it, so no unresolved reference or citation is produced. Printed slips kept literal: no printed word, symbol, hypothesis or number of the counted statement or of its proof was altered. Layout and class: the article's own class is \texttt{amsart} with packages including amssymb, a4wide, amsmath, mathtools, xy, latexsym, pdfsync, xcolor, graphicx, multirow, fontenc, inputenc, array, upgreek; the wrapper is the lane's pinned \texttt{amsart} editorial wrapper at 10pt on A4, which restates the theorem-like environments the retained text needs and adds the article's own macro definitions that the retained text needs (none), with extra wrapper packages (bm, bbm, dsfont, xspace, cancel, mathtools, xcolor). The delivered numbering is the unit's own. Assets: no retained span contains a figure environment, a graphics-inclusion command, a PDF-page inclusion, a table or a TikZ picture; no image file is copied into this package, the package declares no asset dependency and no image file is redistributed with it. Licence: version arXiv:2603.22932v1 of this article is distributed under the Creative Commons Attribution 4.0 International licence, as recorded in the arXiv licence metadata for this exact version and shown on the abstract page https://arxiv.org/abs/2603.22932v1. A full rights notice is delivered at licenses/arxiv-2603.22932-CC-BY-4.0.txt. Characters: every retained excerpt is pure ASCII, so the delivered engine, XeLaTeX with its default Latin Modern text font, reports no missing character.
\begin{equation}
\label{a-antip2}
	\delta_{H}\circ\lambda_{H}=(\lambda_{H}\otimes\lambda_{H})\circ c_{H,H}\circ \delta_{H}{\color{blue} .}
\end{equation}

\begin{gather}\label{compatmodH}
\varphi_{M}^{2}\circ(H\otimes\varphi_{M}^{1})=\varphi_{M}^{1}\circ(\mu_{H}^{2}\otimes\Gamma_{M})\circ(H\otimes c_{H,H}\otimes M)\circ(\delta_{H}\otimes H\otimes M)
\end{gather}

\begin{gather}\label{GammaMdef}
\Gamma_{M}\coloneqq\varphi_{M}^{1}\circ(\lambda_{H}^{1}\otimes\varphi_{M}^{2})\circ(\delta_{H}\otimes M).
\end{gather}

\begin{definition}
	Let $\mathbb{H}$ be a Hopf brace in ${\sf C}$. We will denote by ${}_{\mathbb{H}}{\sf Mod}^{{\sf cc}}$ the full subcategory of ${}_{\mathbb{H}}{\sf Mod}$ whose objects satisfy that $(M,\Gamma_{M})$ belongs to the cocommutativity class of $H_{2}$, that is, for the objects in this subcategory we have that
	\begin{gather}\label{ccclassGammaM}
		(\Gamma_{M}\otimes H)\circ(H\otimes c_{H,M})\circ(\delta_{H}\otimes M)=(\Gamma_{M}\otimes H)\circ(H\otimes c_{H,M})\circ((c_{H,H}\circ\delta_{H})\otimes M).
	\end{gather}
	
	Then, the superscript ${\sf cc}$ in the notation of this subcategory is chosen in reference to the term ``cocommutativity class''. 
	
	If $\mathbb{H}$ is assumed to be cocommutative, then \eqref{ccclassGammaM} holds and, consequently,  ${}_{\mathbb{H}}{\sf Mod}^{{\sf cc}}={}_{\mathbb{H}}{\sf Mod}$.
\end{definition}

\begin{theorem}
	Let $\mathbb{H}=(H_{1},H_{2})$ be a Hopf brace in ${\sf C}$. The category ${}_{\mathbb{H}}{\sf Mod}^{{\sf cc}}$ is monoidal with unit object $(K,\varepsilon_{H}\otimes K,\varepsilon_{H}\otimes K)$ and tensor functor defined as follows:
	\begin{align*}
		\otimes\colon {}_{\mathbb{H}}{\sf Mod}^{{\sf cc}}\times {}_{\mathbb{H}}{\sf Mod}^{{\sf cc}}&\longrightarrow {}_{\mathbb{H}}{\sf Mod}^{{\sf cc}}\\((M,\varphi_{M}^{1},\varphi_{M}^{2}),(N,\varphi_{N}^{1},\varphi_{N}^{2}))&\longmapsto(M\otimes N,\varphi_{M\otimes N}^{1},\varphi_{M\otimes N}^{2}),
	\end{align*}
	where
	\begin{gather}\label{usualtensoraction}
		\varphi_{M\otimes N}^{i}=(\varphi_{M}^{i}\otimes\varphi_{N}^{i})\circ(H\otimes c_{H,M}\otimes N)\circ(\delta_{H}\otimes M\otimes N),\textnormal{ for all }i=1,2.
	\end{gather}
\end{theorem}

\begin{proof}
	It is a classical result in the theory of Hopf algebras that the category of modules over a Hopf algebra $H$ is monoidal, where, if $(M,\varphi_{M})$ and $(N,\varphi_{N})$ are modules over $H$, then the tensor functor is given by $(M\otimes N,\varphi_{M\otimes N})$ with
	\begin{gather*}\label{tensorusualactionhopf}
		\varphi_{M\otimes N}\coloneqq(\varphi_{M}\otimes\varphi_{N})\circ(H\otimes c_{H,M}\otimes N)\circ(\delta_{H}\otimes M\otimes N).\end{gather*}
	Therefore, $(M\otimes N, \varphi_{M\otimes N}^1)$ is a left $H_{1}$-module and $(M\otimes N, \varphi_{M\otimes N}^2)$ is a left $H_{2}$-module.
	
	Let us show that $(M\otimes N,\varphi_{M\otimes N}^{1},\varphi_{M\otimes N}^{2})$ satisfies \eqref{compatmodH}. To do so, note at first that
	\begin{gather}\label{Gammamodtensor}
		\Gamma_{M\otimes N}=(\Gamma_{M}\otimes \Gamma_{N})\circ(H\otimes c_{H,M}\otimes N)\circ(\delta_{H}\otimes M\otimes N).
	\end{gather}
	Indeed,
	\begin{itemize}
		\itemindent=-32pt 
		\item[ ]$\hspace{0.38cm}\Gamma_{M\otimes N}$
		\item[]$=(\varphi_{M}^{1}\otimes\varphi_{N}^{1})\circ(H\otimes c_{H,M}\otimes N)\circ((\delta_{H}\circ\lambda_{H}^{1})\otimes((\varphi_{M}^{2}\otimes\varphi_{N}^{2})\circ(H\otimes c_{H,M}\otimes N)$
		\item[]$\hspace{0.38cm}\circ (\delta_{H}\otimes M\otimes N)))\circ (\delta_{H}\otimes M\otimes N)$ {\footnotesize (by \eqref{GammaMdef} for $(M\otimes N,\varphi_{M\otimes N}^{1},\varphi_{M\otimes N}^{2})$)}
		\item[]$=(\varphi_{M}^{1}\otimes\varphi_{N}^{1})\circ(H\otimes c_{H,M}\otimes N)\circ(((\lambda_{H}^{1}\otimes\lambda_{H}^{1})\circ c_{H,H}\circ\delta_{H})\otimes((\varphi_{M}^{2}\otimes\varphi_{N}^{2})\circ(H\otimes c_{H,M}\otimes N)$
		\item[]$\hspace{0.38cm}\circ (\delta_{H}\otimes M\otimes N)))\circ (\delta_{H}\otimes M\otimes N)$ {\footnotesize (by \eqref{a-antip2} for $H_{1}$)}
		\item[]$=((\varphi_{M}^{1}\circ(\lambda_{H}^{1}\otimes\varphi_{M}^{2}))\otimes(\varphi_{N}^{1}\circ(\lambda_{H}^{1}\otimes\varphi_{N}^{2})))\circ(((H\otimes H\otimes c_{H,M})\circ (H\otimes c_{H,H}\otimes M)$
		\item[]$\hspace{0.38cm}\circ(c_{H,H}\otimes H\otimes M)\circ(H\otimes\delta_{H}\otimes M))\otimes H\otimes N)\circ (((\delta_{H}\otimes c_{H,M})\circ (\delta_{H}\otimes M))\otimes N)$ {\footnotesize (by coassociativity}
		\item[]$\hspace{0.38cm}${\footnotesize of $\delta_{H}$ and the naturality of $c$)}
		\item[]$=(M\otimes (\varphi_{N}^{1}\circ(\lambda_{H}^{1}\otimes \varphi_{N}^{2})))\circ (((\Gamma_{M}\otimes H)\circ(H\otimes c_{H,M})\circ ((c_{H,H}\circ\delta_{H})\otimes M))\otimes H\otimes N)$
		\item[]$\hspace{0.38cm}\circ(H\otimes c_{H,M}\otimes N)\circ(\delta_{H}\otimes M\otimes N)$ {\footnotesize (by naturality of $c$ and \eqref{GammaMdef} for $(M,\varphi_{M}^{1},\varphi_{M}^{2})$)}
		\item[]$=(M\otimes (\varphi_{N}^{1}\circ(\lambda_{H}^{1}\otimes \varphi_{N}^{2})))\circ (((\Gamma_{M}\otimes H)\circ(H\otimes c_{H,M})\circ (\delta_{H}\otimes M))\otimes H\otimes N)$
		\item[]$\hspace{0.38cm}\circ(H\otimes c_{H,M}\otimes N)\circ(\delta_{H}\otimes M\otimes N)$ {\footnotesize (by \eqref{ccclassGammaM})}
		\item[] $=(\Gamma_{M}\otimes\Gamma_{N})\circ(H\otimes c_{H,M}\otimes N)\circ(\delta_{H}\otimes M\otimes N)$ {\footnotesize (by coassociativity of $\delta_{H}$, the naturality of $c$ and}
		\item[]$\hspace{0.38cm}${\footnotesize \eqref{GammaMdef} for $(N,\varphi_{N}^{1},\varphi_{N}^{2})$).}
	\end{itemize}
	
	Then, \eqref{compatmodH} for $(M\otimes N,\varphi_{M\otimes N}^{1},\varphi_{M\otimes N}^{2})$ follows by
	\begin{itemize}
		\itemindent=-32pt 
		\item[ ]$\hspace{0.38cm}\varphi_{M\otimes N}^{1}\circ(\mu_{H}^{2}\otimes \Gamma_{M\otimes N})\circ(H\otimes c_{H,H}\otimes M\otimes N)\circ(\delta_{H}\otimes H\otimes M\otimes N)$
		\item[]$=(\varphi_{M}^{1}\otimes\varphi_{N}^{1})\circ(H\otimes c_{H,M}\otimes N)\circ((\delta_{H}\circ\mu_{H}^{2})\otimes ((\Gamma_{M}\otimes\Gamma_{N})\circ(H\otimes c_{H,M}\otimes N)$
		\item[]$\hspace{0.38cm}\circ(\delta_{H}\otimes M\otimes N)))\circ(H\otimes c_{H,H}\otimes M\otimes N)\circ(\delta_{H}\otimes H\otimes M\otimes N)$ {\footnotesize (by \eqref{usualtensoraction} and \eqref{Gammamodtensor})}
		\item[]$=((\varphi_{M}^{1}\circ (\mu_{H}^{2}\otimes \Gamma_{M}))\otimes(\varphi_{N}^{1}\circ(\mu_{H}^{2}\otimes\Gamma_{N})))\circ (H\otimes((H\otimes H\otimes c_{H,M}\otimes H\otimes H)\circ(((H\otimes c_{H,H})$
		\item[]$\hspace{0.38cm}\circ(c_{H,H}\otimes H)\circ(H\otimes c_{H,H})\circ (\delta_{H}\otimes H))\otimes((c_{H,M}\otimes H)\circ (H\otimes c_{H,M})\circ (c_{H,H}\otimes M)))$
		\item[]$\hspace{0.38cm}\circ (H\otimes H\otimes c_{H,H}\otimes H\otimes M))\otimes N)\circ(((\delta_{H}\otimes H)\circ\delta_{H})\otimes\delta_{H}\otimes M\otimes N)$ {\footnotesize (by naturality of $c$, the}
		\item[]$\hspace{0.38cm}${\footnotesize symmetric character of ${\sf C}$ and coassociativity of $\delta_{H}$)}
		\item[]$=((\varphi_{M}^{1}\circ (\mu_{H}^{2}\otimes M))\otimes(\varphi_{N}^{1}\circ(\mu_{H}^{2}\otimes\Gamma_{N})))\circ(H\otimes((H\otimes ((\Gamma_{M}\otimes H)\circ(H\otimes c_{H,M})$
		\item[]$\hspace{0.38cm}\circ((c_{H,H}\circ\delta_{H})\otimes M))\otimes H\otimes H)\circ (c_{H,H}\otimes ((c_{H,M}\otimes H)\circ (H\otimes c_{H,M})\circ (c_{H,H}\otimes M)))$
		\item[]$\hspace{0.38cm}\circ(H\otimes c_{H,H}\otimes H\otimes M))\otimes N)\circ(((\delta_{H}\otimes H)\circ\delta_{H})\otimes\delta_{H}\otimes M\otimes N)$ {\footnotesize (by naturality of $c$)}
		\item[]$=((\varphi_{M}^{1}\circ (\mu_{H}^{2}\otimes M))\otimes(\varphi_{N}^{1}\circ(\mu_{H}^{2}\otimes\Gamma_{N})))\circ(H\otimes((H\otimes ((\Gamma_{M}\otimes H)\circ(H\otimes c_{H,M})$
		\item[]$\hspace{0.38cm}\circ(\delta_{H}\otimes M))\otimes H\otimes H)\circ (c_{H,H}\otimes ((c_{H,M}\otimes H)\circ (H\otimes c_{H,M})\circ (c_{H,H}\otimes M)))$
		\item[]$\hspace{0.38cm}\circ(H\otimes c_{H,H}\otimes H\otimes M))\otimes N)\circ(((\delta_{H}\otimes H)\circ\delta_{H})\otimes\delta_{H}\otimes M\otimes N)$ {\footnotesize (by \eqref{ccclassGammaM})}
		\item[]$=((\varphi_{M}^{1}\circ(\mu_{H}^{2}\otimes\Gamma_{M})\circ(H\otimes c_{H,H}\otimes M)\circ(\delta_{H}\otimes H\otimes M))\otimes(\varphi_{N}^{1}\circ (\mu_{H}^{2}\otimes\Gamma_{N})$
		\item[]$\hspace{0.38cm}\circ (H\otimes c_{H,H}\otimes N)))\circ (H\otimes((H\otimes c_{H,M}\otimes H)\circ (c_{H,H}\otimes c_{H,M})\circ (H\otimes c_{H,H}\otimes M)$
		\item[]$\hspace{0.38cm}\circ(\delta_{H}\otimes H\otimes M))\otimes H\otimes N)\circ (\delta_{H}\otimes((H\otimes c_{H,M})\circ(\delta_{H}\otimes M))\otimes N)$ {\footnotesize (by naturality of $c$ and}
		\item[]$\hspace{0.38cm}${\footnotesize coassociativity of $\delta_{H}$)}
		\item[]$=((\varphi_{M}^{2}\circ(H\otimes\varphi_{M}^{1}))\otimes(\varphi_{N}^{1}\circ(\mu_{H}^{2}\otimes \Gamma_{N})\circ(H\otimes c_{H,H}\otimes N)\circ(\delta_{H}\otimes H\otimes N)))$
		\item[]$\hspace{0.38cm}\circ (H\otimes((H\otimes c_{H,M}\otimes H)\circ(c_{H,H}\otimes c_{H,M})) \otimes N)\circ(\delta_{H}\otimes\delta_{H}\otimes M\otimes N)$ {\footnotesize (by naturality of $c$ and}
		\item[]$\hspace{0.38cm}${\footnotesize \eqref{compatmodH} for $(M,\varphi_{M}^{1},\varphi_{M}^{2})$)}
		\item[]$=(\varphi_{M}^{2}\otimes (\varphi_{N}^{2}\circ(H\otimes \varphi_{N}^{1})))\circ (H\otimes((\varphi_{M}^{1}\otimes H)\circ (H\otimes c_{H,M})\circ (c_{H,H}\otimes M))\otimes H\otimes N)$
		\item[]$\hspace{0.38cm}\circ(\delta_{H}\otimes((H\otimes c_{H,M})\circ(\delta_{H}\otimes M))\otimes N)$ {\footnotesize (by \eqref{compatmodH} for $(N,\varphi_{N}^{1},\varphi_{N}^{2})$)} 
		\item[]$=\varphi_{M\otimes N}^{2}\circ(H\otimes\varphi_{M\otimes N}^{1})$ {\footnotesize (by naturality of $c$ and \eqref{usualtensoraction})}.
	\end{itemize}
	
	To conclude, it only remains to show that \eqref{ccclassGammaM} holds for $(M\otimes N,\varphi_{M\otimes N}^{1},\varphi_{M\otimes N}^{2})$.
	Indeed,
	\begin{itemize}
		\itemindent=-32pt 
		\item[ ]$\hspace{0.38cm}\Gamma_{M\otimes N}\circ(H\otimes c_{H,M\otimes N})\circ((c_{H,H}\circ\delta_{H})\otimes M\otimes N)$
		\item[]$=(((\Gamma_{M}\otimes\Gamma_{N})\circ (H\otimes c_{H,M}\otimes N)\circ (\delta_{H}\otimes M\otimes N))\otimes H)\circ (H\otimes((M\otimes c_{H,N})\circ (c_{H,M}\otimes N)))$
		\item[]$\hspace{0.38cm}\circ((c_{H,H}\circ\delta_{H})\otimes M\otimes N)$ {\footnotesize (by \eqref{Gammamodtensor})}
		\item[]$=(M\otimes ((\Gamma_{N}\otimes H)\circ(H\otimes c_{H,N})\circ (c_{H,H}\otimes N)))\circ(((\Gamma_{M}\otimes H)\circ(H\otimes c_{H,M})\circ((c_{H,H}\circ\delta_{H})$
		\item[]$\hspace{0.38cm}\otimes M))\otimes H\otimes N)\circ (((H\otimes c_{H,M})\circ(\delta_{H}\otimes M))\otimes N)$ {\footnotesize (by naturality of $c$ and coassociativity of $\delta_{H}$)}
		\item[]$=(M\otimes ((\Gamma_{N}\otimes H)\circ(H\otimes c_{H,N})\circ (c_{H,H}\otimes N)))\circ(((\Gamma_{M}\otimes H)\circ(H\otimes c_{H,M})\circ(\delta_{H}\otimes M))$
		\item[]$\hspace{0.38cm}\otimes H\otimes N)\circ (((H\otimes c_{H,M})\circ(\delta_{H}\otimes M))\otimes N)$ {\footnotesize (by \eqref{ccclassGammaM} for $(M,\Gamma_{M})$)}
		\item[]$=(\Gamma_{M}\otimes((\Gamma_{N}\otimes H)\circ (H\otimes c_{H,N})\circ ((c_{H,H}\circ\delta_{H})\otimes N)))\circ(H\otimes c_{H,M}\otimes N)\circ(\delta_{H}\otimes M\otimes N)$ 
		\item[]$\hspace{0.38cm}${\footnotesize (by coassociativity of $\delta_{H}$ and naturality of $c$)}
		\item[]$=(\Gamma_{M}\otimes((\Gamma_{N}\otimes H)\circ (H\otimes c_{H,N})\circ (\delta_{H}\otimes N)))\circ(H\otimes c_{H,M}\otimes N)\circ(\delta_{H}\otimes M\otimes N)$ {\footnotesize (by \eqref{ccclassGammaM}}
		\item[]$\hspace{0.38cm}${\footnotesize for $(N,\Gamma_{N})$)}
		\item[]$=\Gamma_{M\otimes N}\circ(H\otimes c_{H,M\otimes N})\circ(\delta_{H}\otimes M\otimes N)$ {\footnotesize (by naturality of $c$ and coassociativity of $\delta_{H}$)}.\qedhere
	\end{itemize}
\end{proof}

\end{document}
